\chapter{The dataset}
\label{ch:dataset}

\section{Source}
\label{sec:source}

The flights come from the \emph{Coupe F\'ed\'erale de Distance} (CFD), the
cross-country competition of the French free-flight federation (FFVL). Pilots declare
their flights online, and each entry carries the flight's metadata and, in most
cases, its GPS track. Paragliders and hang gliders have separate archives; a larger
share of hang-glider entries has no track.

A season runs from 1~September to 31~August. We use the paraglider seasons from
1999--2000 and the hang-glider seasons from 2001--2002, up to 2025--2026. For each
season we download the list of declared flights, kept verbatim as provenance, and
every track it links. A saved list is never replaced, so each one is a snapshot of
the day it was downloaded: the 2025--2026 lists were saved around the end of June
2026 and miss the rest of that season.

\section{Tracks and metadata}
\label{sec:tracks}

Tracks are IGC files~\cite{fai_igc_spec}, plain-text records of the flight. Each
\texttt{B} record is one \emph{fix}: UTC time, latitude and longitude, a validity flag,
and two altitudes, the pressure altitude from the barometer and the GNSS altitude from
the satellite fix. The sampling interval depends on the recorder and need not be
constant within a flight. Some recorders log no pressure altitude.

The metadata give the date, the pilot, the take-off and landing sites, the wing
model and class, the scored distance and the scoring category of the declared route
(free distance, out-and-return, flat or FAI triangle and a few others). The category
describes the declared route only. A few paraglider entries have an incomplete date,
such as \texttt{0000-00-00}, which excludes them from any comparison by date.

\section{Glider categories}
\label{sec:glider}

Wing design changes speed, glide performance and turning geometry, so paragliders and
hang gliders are analysed separately~\cite{vilpellet2026}. A paraglider is a flexible
canopy; a hang glider has a framed wing, flexible or rigid.

The wing class follows a different system in each discipline. Paraglider classes are
those of the EN~926-2 standard, A to D, written with their equivalent in the older LTF
scale (\enquote{C ou 2} reads \enquote{C or 2}), plus the CIVL Competition Class and a
few other labels, such as tandem or uncertified wings. EN classes rank passive
safety and the skill needed to recover from a collapse, and say nothing direct about speed or glide
ratio~\cite{fpren926_2_2012}. Hang-glider classes are those of the
FAI~\cite{fai_sc7_2024}: Class~1, with its Sport sub-class, for wings controlled
by weight shift alone, and Classes~2 and~5 for rigid wings with movable control
surfaces.
